\documentclass[xcolor={dvipsnames}]{beamer}

\usepackage{amsmath,amssymb,amsfonts,amsthm,mathrsfs,graphicx,setspace,calligra}
\usepackage[dvipsnames]{xcolor}
\usepackage{hyperref}

\usetheme{Madrid}
\definecolor{myteal}{cmyk}{0.5,0,0.15,0}
\definecolor{deepblue}{RGB}{0,36,96}
\usecolortheme[named=myteal]{structure}
\usefonttheme{professionalfonts}
\useoutertheme{infolines}

\setbeamersize{text margin left=0.45cm,text margin right=0.45cm}
\setbeamertemplate{navigation symbols}{}
\setbeamerfont{block title}{size=\normalsize}
\setbeamerfont{block body}{size=\small}
\setbeamercolor{frametitle}{fg=deepblue,bg=myteal}
\setbeamercolor{title}{fg=deepblue}
\setbeamercolor{date in head/foot}{fg=deepblue,bg=myteal}
\setbeamercovered{transparent}
\newcommand{\focus}[1]{%
  \begingroup\setlength{\fboxsep}{1.5pt}%
  \colorbox{myteal!18}{\textcolor{deepblue}{\bfseries #1}}%
  \endgroup}

\theoremstyle{definition}
\newtheorem{df}{Definition}
\newtheorem{eg}[df]{Example}
\newtheorem{ques}[df]{Question}
\newtheorem{re}[df]{Remark}
\theoremstyle{plain}
\newtheorem{thm}[df]{Theorem}
\newtheorem{pp}[df]{Proposition}
\newtheorem{co}[df]{Corollary}

\newcommand{\MS}{\mathcal{S}}
\newcommand{\MA}{\mathcal{A}}
\newcommand{\MB}{\mathcal{B}}
\newcommand{\MH}{\mathcal{H}}
\newcommand{\MJ}{\mathcal{J}}
\newcommand{\MG}{\mathscr{G}}
\newcommand{\Rbb}{\mathbb{R}}
\newcommand{\Zbb}{\mathbb{Z}}
\newcommand{\Cbb}{\mathbb{C}}
\newcommand{\cl}{\mathrm{cl}}
\newcommand{\op}{\mathrm{op}}
\newcommand{\loc}{\mathrm{loc}}
\newcommand{\Rep}{\mathrm{Rep}}
\newcommand{\End}{\mathrm{End}}
\newcommand{\Aut}{\mathrm{Aut}}
\newcommand{\Diff}{\mathrm{Diff}}
\newcommand{\Cf}{\mathrm{Cf}}
\newcommand{\Mob}{\mathrm{Mob}}
\newcommand{\Vir}{\mathrm{Vir}}
\newcommand{\Ad}{\mathrm{Ad}}
\newcommand{\Hom}{\mathrm{Hom}}
\newcommand{\id}{\mathrm{id}}
\newcommand{\rint}{\mathrm{Int}}
\newcommand{\Rcl}{\mathbb R^{1,1}_{\cl}}
\newcommand{\Rop}{\mathbb R^{1,1}_{\op}}
\newcommand{\Bcl}{\MB_{\cl}}
\newcommand{\Hcl}{\MH_{\cl}}
\newcommand{\Hop}{\MH_{\op}}
\newcommand{\Hopij}[2]{\MH_{\op(#1,#2)}}
\newcommand{\Bopi}[1]{\MB_{\op(#1)}}
\newcommand{\Jt}{\widetilde{\MJ}}
\newcommand{\Gc}{\MG_c}
\newcommand{\Hsqrt}[1]{\MH_{\sqrt{#1}}}
\newcommand{\bA}{\boxtimes_{\MA}}
\newcommand{\bAtwo}{\boxtimes_{\MA^{\otimes 2}}}

\newcommand{\ClosedOpenSpaces}{%
  \includegraphics[width=\linewidth]{figures/crops/fig_spaces.png}}
\newcommand{\BoundaryHaagPicture}{%
  \includegraphics[width=.4\linewidth]{figures/crops/fig1d.pdf}}
\newcommand{\ClosedCylinderPicture}{%
  \includegraphics[height=.22\textheight]{figures/crops/fig_closed_cyl_O.png}}
\newcommand{\SpacelikePicture}{%
  \includegraphics[width=.82\textwidth]{figures/crops/fig_spacelike_locality.png}}
\newcommand{\ClosedHaagPicture}{%
  \includegraphics[width=.52\linewidth]{figures/crops/fig_closed_multi.png}}
\newcommand{\BulkBoundaryPicture}{%
  \includegraphics[width=.78\linewidth]{figures/crops/fig_bulk_boundary.png}}
\newcommand{\PartFrame}[2]{%
\begin{frame}
\vfill
\begin{center}
{\large\textcolor{myteal!70!black}{Part #1}}\\[.35cm]
{\LARGE\textcolor{deepblue}{#2}}
\end{center}
\vfill
\end{frame}}

\title[Open/closed CFT]{Minkowskian open/closed conformal field theory\\possibly without vacuum: the Cardy case}
\author[Bin Gui]{\texorpdfstring{Bin Gui\\Tsinghua University}{Bin Gui, Tsinghua University}}
\date[arXiv:2606.30928]{ZMP Colloquium, Universit\"at Hamburg\\July 2026\\arXiv:2606.30928 with the same title}

\begin{document}
\begin{spacing}{1.04}

\frame{\titlepage}

\begin{frame}[t]{Outline}
\vspace{.45cm}
\begin{center}
\begin{minipage}{.86\textwidth}
\renewcommand{\arraystretch}{1.55}
\begin{tabular}{@{}p{.9cm}@{\hspace{.35cm}}p{8.95cm}@{}}
{\Large\textcolor{myteal!70!black}{I}} &
{\large\textcolor{deepblue}{Geometric setup}}\\[-.03cm]
\multicolumn{2}{@{}c@{}}{\textcolor{myteal!45}{\rule{\linewidth}{.35pt}}}\\[-.05cm]
{\Large\textcolor{myteal!70!black}{II}} &
{\large\textcolor{deepblue}{Boundary fields and open string state spaces}}\\[-.03cm]
\multicolumn{2}{@{}c@{}}{\textcolor{myteal!45}{\rule{\linewidth}{.35pt}}}\\[-.05cm]
{\Large\textcolor{myteal!70!black}{III}} &
{\large\textcolor{deepblue}{Bulk fields and closed string state space}}\\[-.03cm]
\multicolumn{2}{@{}c@{}}{\textcolor{myteal!45}{\rule{\linewidth}{.35pt}}}\\[-.05cm]
{\Large\textcolor{myteal!70!black}{IV}} &
{\large\textcolor{deepblue}{Bulk fields acting on open string state spaces}}
\end{tabular}
\end{minipage}
\end{center}
\end{frame}

\PartFrame{I}{Geometric Setup}

\begin{frame}[t]{Closed and Open Strings}
\footnotesize
\begin{block}{Closed and open strings}
2D CFT studies the propagation and interaction of closed strings
\[
  S^1\simeq \Rbb/2\pi\Zbb
\]
and open strings
\[
  [0,\pi].
\]
The corresponding Minkowski space-times are
\[
  \Rcl=S^1\times\Rbb,\qquad
  \Rop=[0,\pi]\times\Rbb .
\]
\end{block}

\begin{columns}[T,onlytextwidth]
\column{.5\textwidth}
\begin{align*}
  &(x,t)\in \Rcl\ \hbox{or}\ \Rop,\\
  &(ds)^2=(dt)^2-(dx)^2,\\
  &u^+=t+x,\qquad u^-=t-x,\\
  &(ds)^2=du^+du^- .
\end{align*}
\column{.47\textwidth}
\begin{center}
\ClosedOpenSpaces
\end{center}
\end{columns}
\end{frame}

\begin{frame}[t]{Conformal Diffeomorphisms}
\footnotesize
Let $\Diff_+(S^1)$ be the group of orientation-preserving diffeomorphisms of $S^1$, and let
\[
  \MG=\widetilde{\Diff_+(S^1)}
\]
be its universal cover.
Here
\[
  \Cf^+(\Rcl),\ \Cf^+(\Rop)
\]
denote orientation-preserving, time-direction-preserving conformal diffeomorphism groups.
Elements $F\in \MG$ can be viewed as smooth functions $F:\Rbb\to\Rbb$ such that
\[
  F(x+2\pi)=F(x)+2\pi,\qquad F'(x)>0 .
\]

In light-cone coordinates, the relevant actions are
\[
  (g_1,g_2)(u^+,u^-)=(g_1u^+,g_2u^-)
  \quad\hbox{on }\Rcl,
\]
\[
  g(u^+,u^-)=(gu^+,gu^-)
  \quad\hbox{on }\Rop.
\]

Both actions are well-defined.
This allows us to classify $\Cf^+(\Rcl)$ and $\Cf^+(\Rop)$.
\end{frame}

\begin{frame}[t]{$\Cf^+(\Rcl)$}
\small
The action of $\MG\times \MG$ on $\Rcl$ gives the covering exact sequence
\[
  1\longrightarrow \Zbb\longrightarrow
  \MG\times \MG\longrightarrow \Cf^+(\Rcl)\longrightarrow 1.
\]

\begin{block}{Kernel}
The subgroup $\Zbb$ is generated by
\[
  \bigl(\rho(2\pi),\rho(-2\pi)\bigr)\in \MG\times \MG,
\]
where
\[
  \rho:\Rbb\longrightarrow \MG
\]
is the one-parameter rotation subgroup.
\end{block}

\end{frame}

\begin{frame}[t]{$\Cf^+(\Rop)$}
\small
Let $\MG$ act on $\Rop$ by
\[
  g(u^+,u^-)=(gu^+,gu^-).
\]

\begin{block}{Classification}
The action is well-defined and gives
\[
  \MG\simeq \Cf^+(\Rop).
\]
\end{block}

\begin{itemize}
  \item Closed CFT uses $\MG\times \MG$.
  \item Open CFT uses $\MG$.
  \item With central charge $c$, we use the central extension $\Gc$.
\end{itemize}
\end{frame}

\PartFrame{II}{Boundary Fields and Open String State Spaces}

\begin{frame}[t]{Adapting Fuchs--Schweigert--Runkel}
\footnotesize
Adapting Fuchs--Schweigert--Runkel to the operator algebraic approach of CFT, we expect a 2D open/closed CFT with central charge $c$ as follows.

\begin{itemize}
  \item Boundary conditions are labeled by indices
  \[
    i,j,k,\ldots .
  \]
  \item The open-string state space with left and right boundary conditions $i,j$ is
  \[
    \Hopij{i}{j}.
  \]
  \item Boundary operators preserving a boundary condition $i$ form a generally non-local net
  \[
    \Bopi{i}:\ \widetilde I\in\Jt\longmapsto
    \Bopi{i}(\widetilde I)\subset \mathfrak L(\Hopij{i}{i}).
  \]
where $\Bopi{i}(\widetilde I)$ is a von Neumann algebra (i.e. a $*$-subalgebra closed under strong operator topology).
\end{itemize}

\begin{block}{Arg-valued intervals}
\[
  \Jt=\{\widetilde I:\ \widetilde I
  \hbox{ is a component over } I\subset S^1\}
\]
\end{block}
\end{frame}

\begin{frame}[t]{$\Jt$}
\small
Let
\[
  \MJ=\{I\subset S^1:\ I\hbox{ is non-empty, open, non-dense, and connected}\}.
\]
Let
\[
  \Jt=\{\widetilde I:\ \widetilde I
  \hbox{ is a component of } I
  \hbox{ in } \Rbb\to\Rbb/2\pi\Zbb=S^1\}.
\]
Equivalently,
\[
  \Jt=\{\hbox{open intervals in }\Rbb\hbox{ of length }0<|\widetilde I|<2\pi\}.
\]

\begin{block}{Non-local boundary net}
For a boundary condition $i$,
\[
  \Bopi{i}:\widetilde I\in\Jt\longmapsto
  \Bopi{i}(\widetilde I)\subset \mathfrak L(\Hopij{i}{i})
\]
where $ \Bopi{i}(\widetilde I)$ is a von Neumann algebra.
\end{block}
\end{frame}

\begin{frame}[t]{Conformal Covariance}
\small
The boundary net is conformally covariant.  Let $\Gc$ be the central extension of $\MG$ with central charge $c$.

\begin{block}{Covariance}
There is a strongly continuous unitary representation of $\Gc$ on $\Hopij{i}{i}$, indeed on $\Hopij{i}{j}$, such that
\[
  g\,\Bopi{i}(\widetilde I)\,g^{-1}
  =
  \Bopi{i}(g\widetilde I),
  \qquad \widetilde I\in\Jt .
\]
\end{block}

The same convention applies to the right action of $\Bopi{i}$.
\end{frame}

\begin{frame}[t]{Left and Right Actions}
\small
On the open-string state space $\Hopij{i}{j}$, we have a left action of $\Bopi{i}$ and a right action of $\Bopi{j}$:
\[
  \Bopi{i}(\widetilde I)
  \curvearrowright
  \Hopij{i}{j}
  \curvearrowleft
  \Bopi{j}(\widetilde J),
\]
where the first one is the left action and the second one is the right action.

\begin{block}{Boundary-boundary locality}
If $\widetilde J$ is clockwise to $\widetilde I$, i.e. $0<\arg z-\arg\zeta<2\pi$ for each $z\in\widetilde I$ and $\zeta\in \widetilde J$, then
\[
  \bigl[
    \Bopi{i}(\widetilde I),
    \Bopi{j}(\widetilde J)
  \bigr]=0.
\]
\end{block}


\end{frame}

\begin{frame}[t]{Boundary-Boundary Haag Duality}
\small
\begin{columns}[T,onlytextwidth]
\column{.48\textwidth}
\begin{center}
\BoundaryHaagPicture
\end{center}

\column{.5\textwidth}
For $\Hopij{i}{j}$, the nets $\Bopi{i}$ and $\Bopi{j}$ act on the left and right boundaries.

\begin{block}{Boundary-boundary Haag duality}
The following von Neumann algebras are commutants of each other:
\[
  \Bopi{i}(\widetilde I),
  \qquad
  \Bopi{j}(\widetilde I').
\]
This is called the \focus{boundary-boundary Haag duality}.
\end{block}
\end{columns}

\vspace{.1cm}
Here, $\widetilde I'$ is the largest element in $\widetilde{\mathcal J}$ clockwise to $\widetilde I$. In particular, if $\widetilde J$ is clockwise to $\widetilde I$, then
\[
  \bigl[
  \Bopi{i}(\widetilde I),
  \Bopi{j}(\widetilde J)
  \bigr]=0 .
\]
\end{frame}

\begin{frame}[t]{Morita Equivalent Boundary Conditions}
\small
That $\Bopi{i}$ and $\Bopi{j}$ are commutants of each other on $\Hopij{i}{j}$ implies that $\Bopi{i}$ and $\Bopi{j}$ are \focus{Morita equivalent}.

This would imply that they produce the same closed CFT. This still needs to be proved in the operator algebra framework.

\begin{block}{Conjugation}
\[
  \overline{\Hopij{i}{j}}=\Hopij{j}{i}.
\]
\end{block}

\begin{block}{Composition of open strings}
\[
  \Hopij{i}{j}\boxtimes_{\Bopi{j}}\Hopij{j}{k}
  \simeq
  \Hopij{i}{k}.
\]
Here $\boxtimes$ denotes the Connes fusion.
\end{block}
\end{frame}

\begin{frame}[t]{Cardy-Type CFT}
\small
We say that the CFT is of Cardy type if there is a boundary condition $0$ such that
\[
  \bigl(\Bopi{0},\Hopij{0}{0}\bigr)
\]
is a local conformal net
\[
  (\MA,\MH_0).
\]
This means $(\MA,\MH_0)$ is a non-local net satisfying the following conditions:

\begin{enumerate}
  \item single-valuedness on $S^1$;
  \item locality:
  \[
    I\cap J=\emptyset\quad\Rightarrow\quad
    [\MA(I),\MA(J)]=0;
  \]
  \item a vacuum vector $\Omega\in\MH_0$ fixed by the Mobius group;
  \item and some other properties.
\end{enumerate}
\end{frame}

\begin{frame}[t]{The Local Conformal Net $(\MA,\MH_0)$}
\scriptsize
\vspace{-.15cm}
The pair $(\MA,\MH_0)$ is a non-local net satisfying the following extra conditions.

\begin{block}{Single-valuedness}
If $\widetilde I_1,\widetilde I_2\in\Jt$ project to the same interval $I\in\MJ$, then
\[
  \MA(\widetilde I_1)=\MA(\widetilde I_2).
\]
Thus we write
\[
  I\in\MJ\longmapsto \MA(I)\subset\mathfrak L(\MH_0).
\]
\end{block}

\begin{block}{Locality and vacuum}
\[
  I\cap J=\emptyset \quad\Rightarrow\quad [\MA(I),\MA(J)]=0.
\]
There is a vacuum vector $\Omega\in\MH_0$ fixed by the action of
\[
  \Mob\subset \MG,
\]
the subgroup of linear fractional transformations preserving $S^1$.
\end{block}
\end{frame}

\begin{frame}[t]{The Cardy-type CFT}
\small
Now restrict to the Cardy-type CFT.  Set
\[
  \MH_i:=\Hopij{i}{0}.
\]
Then $\MH_i$ is a right $\MA$-module. %Thus boundary conditions are labeled by right $\MA$-modules, also called solitonic $\MA$-modules; for this talk we take local $\MA$-modules.

\begin{block}{Open-string state space}
Since one expects
\[
  \Hopij{i}{j}
  =
  \Hopij{i}{0}\boxtimes_{\Bopi{0}}\Hopij{0}{j}
  =
  \Hopij{i}{0}\boxtimes_{\Bopi{0}}\overline{\Hopij{j}{0}},
\]
we have
\[
  \Hopij{i}{j}
  =
  \MH_i\bA\overline{\MH_j}.
\]
\end{block}
\end{frame}

\begin{frame}[t]{Boundary Conditions Labeled by $\MA$-Modules}
\small
The boundary conditions in the Cardy case are labeled by right $\MA$-modules, equivalently solitonic $\MA$-modules.

\begin{block}{Local modules}
For simplicity, in this talk we restrict to local $\MA$-modules:
\[
  \hbox{left action}=\hbox{right action}.
\]
We simply call them $\MA$-modules.
\end{block}

\begin{block}{Example}
If $\MA$ is the Virasoro conformal net, then $\MA$-modules correspond to unitary representations of the Virasoro algebra.
\end{block}
\end{frame}

\PartFrame{III}{Bulk Fields and Closed String State Space}

\begin{frame}[t]{Closed CFT}
\footnotesize
For each local $\MA$-module $\MH_i$, there is a canonical $\sigma$-twisted $\MA^{\otimes 2}$-module structure on the same Hilbert space, where
\[
  \sigma\in\Aut(\MA^{\otimes 2})
\]
is the permutation of the two tensor factors.

\begin{block}{Longo--Xu}
We denote this twisted module by
\[
  \Hsqrt{i}.
\]
This is the Longo--Xu result (04); a vertex-algebra version appeared in Barron--Dong--Mason (02).
\end{block}
\end{frame}

\begin{frame}[t]{The State Space for the Closed CFT}
\footnotesize
The state space for the closed CFT is
\[
  \Hcl
  =
  \Hsqrt{0}\bAtwo\overline{\Hsqrt{0}},
\]
as an untwisted $\MA^{\otimes 2}$-module.

\begin{columns}[T,onlytextwidth]
\begin{column}{.70\textwidth}
\begin{block}{Bulk net}
We have a net of von Neumann algebras
\[
  \Bcl:\ O\longmapsto \Bcl(O)\subset\mathfrak L(\Hcl),
\]
where $O$ is a double cone in $\Rcl$.
\end{block}
\end{column}
\begin{column}{.27\textwidth}
\begin{center}
\vspace{.15cm}
\ClosedCylinderPicture
\end{center}
\end{column}
\end{columns}
\vspace{.25cm}

The net satisfies properties such as
\[
  \Gc\times\Gc\hbox{-covariance},\qquad
  \hbox{locality},\qquad
  \hbox{\focus{bulk-bulk Haag duality}}.
\]
\end{frame}

\begin{frame}[t]{Bulk-Bulk Haag Duality}
\footnotesize
\begin{columns}[T,onlytextwidth]
\column{.43\textwidth}
\begin{center}
\includegraphics[width=.88\linewidth]{figures/crops/fig_bulk_haag.png}
\end{center}

\column{.55\textwidth}
\begin{block}{Single double cone}
If $O'$ is the spacelike complement of $O$, then
\[
  \Bcl(O)'=\Bcl(O').
\]
\end{block}
\end{columns}

\vspace{.08cm}
\begin{block}{Multi double-cones}
If $O_1,\ldots,O_n$ are mutually spacelike separated and
\[
  (O_1\cup\cdots\cup O_n)'=D_1\cup\cdots\cup D_n,
\]
then
\[
  \bigl(\Bcl(O_1)\vee\cdots\vee\Bcl(O_n)\bigr)'
  =
  \Bcl(D_1)\vee\cdots\vee\Bcl(D_n).
\]
\end{block}
\end{frame}

\begin{frame}[t]{Two and Multi Double-Cones}
\footnotesize
\begin{columns}[T,onlytextwidth]
\column{.45\textwidth}
\begin{center}
\vspace{.25cm}
\includegraphics[width=\linewidth]{figures/crops/fig1c.pdf}
\end{center}

\column{.53\textwidth}
\begin{block}{Two double-cones version}
In the configuration shown on the left,
\[
  (O_1\cup O_2)'=D_1\cup D_2,
\]
and
\[
\begin{aligned}
  &\bigl(\Bcl(O_1)\vee\Bcl(O_2)\bigr)'\\
  &\qquad =\Bcl(D_1)\vee\Bcl(D_2).
\end{aligned}
\]
\end{block}
\end{columns}

\vspace{.15cm}
The same statement holds for multi double-cones.
\end{frame}

\begin{frame}[t]{Modular Invariance}
\small
\begin{columns}[T,onlytextwidth]
\column{.46\textwidth}
\begin{center}
\includegraphics[width=\linewidth]{figures/crops/fig1c.pdf}
\end{center}

\column{.50\textwidth}
In Euclidean CFT, this version of Haag duality corresponds to \focus{modular invariance}.
\end{columns}

\vspace{.45cm}
Multi double-cone version also holds.
\end{frame}

\PartFrame{IV}{Bulk Fields Acting on Open String State Spaces}

\begin{frame}[t]{$\Bcl$ Acts on $\Hopij{i}{j}$}
\small
The closed net restricts to the open-string state space
\[
  \Hopij{i}{j}=\MH_i\bA\overline{\MH_j}.
\]
Thus, for each double cone $O\subset \rint\Rop=(0,\pi)\times\Rbb$, we get an action
\[
  \Bcl(O)\curvearrowright\Hopij{i}{j}.
\]

The action satisfies, among other things, the \focus{bulk-boundary Haag duality}.
\end{frame}

\begin{frame}[t]{Bulk-Boundary Haag Duality}
\footnotesize
\begin{columns}[T,onlytextwidth]
\column{.43\textwidth}
\begin{center}
\includegraphics[width=.72\linewidth]{figures/crops/fig1e.pdf}
\end{center}

\column{.55\textwidth}
Let $\widetilde E_+$ and $\widetilde E_-$ be the two boundary intervals in the spacelike complement of the bulk double cone $O$, as in the picture.

\begin{block}{Bulk-boundary Haag duality}
\scriptsize
\[
  \Bcl(O)
  \quad\hbox{and}\quad
  \Bopi{i}(\widetilde E_+)\vee\Bopi{j}(\widetilde E_-)
\]
are commutants of each other.
\end{block}
\end{columns}

\vspace{.1cm}
This corresponds to the \focus{Cardy consistency condition} in Euclidean CFT.
\end{frame}

\begin{frame}[t]{Multi Double-Cones in the Strip}
\small
\begin{center}
\includegraphics[width=.38\linewidth]{figures/crops/fig1f.pdf}
\end{center}

\vspace{.2cm}
The multi double-cone version holds.
\end{frame}

\begin{frame}[t]{Theorem}
\small
\vspace{.35cm}
\begin{center}
\begin{minipage}{.92\textwidth}
\begin{block}{Theorem}
For any conformal net $\MA$, there exists the open/closed CFT satisfying the previously mentioned properties, whose boundary conditions are labeled by local $\MA$-modules.
\end{block}
\end{minipage}
\end{center}

\vspace{.45cm}
\begin{minipage}{.92\textwidth}
\footnotesize
\textcolor{deepblue}{\textbf{Expectations.}}\quad
We expect the construction to generalize to:
\begin{itemize}
  \item non-local $\MA$-modules;
  \item or even general non-local nets.
\end{itemize}

\vspace{.05cm}
We expect that the open/closed CFT for $\MA=$ Virasoro conformal net of $c\geq 25$ is the Minkowskian open/closed Liouville CFT.

\smallskip
More generally, Toda CFTs are expected to be Cardy-type CFTs associated with principal $\mathcal W(\mathfrak g)$ conformal nets at Toda central charge.
\end{minipage}
\end{frame}

\begin{frame}[t]{Three Haag Dualities}
\footnotesize
\vspace{-.05cm}
\begin{center}
\renewcommand{\arraystretch}{1.18}
\begin{tabular}{|l|c|l|}
\hline
\textbf{Haag duality} & \textbf{Spacetime} & \textbf{Interpretation}\\
\hline
Bulk-bulk & $\Rcl$ & \begin{tabular}[c]{@{}l@{}}Modular invariance\\in Euclidean spacetimes\end{tabular}\\
\hline
Boundary-boundary & $\Rop$ & Morita equivalence of boundary nets\\
\hline
Bulk-boundary & $\Rop$ & \begin{tabular}[c]{@{}l@{}}Cardy consistency condition\\in Euclidean spacetimes\end{tabular}\\
\hline
\end{tabular}
\end{center}

\vspace{.12cm}
\begin{columns}[T,onlytextwidth]
\column{.34\textwidth}
\centering
\includegraphics[height=.22\textheight,width=.98\linewidth,keepaspectratio]{figures/crops/fig1c.pdf}\\[-.05cm]
\scriptsize Bulk-bulk

\column{.31\textwidth}
\centering
\includegraphics[height=.22\textheight,width=.98\linewidth,keepaspectratio]{figures/crops/fig1d.pdf}\\[-.05cm]
\scriptsize Boundary-boundary

\column{.31\textwidth}
\centering
\includegraphics[height=.22\textheight,width=.98\linewidth,keepaspectratio]{figures/crops/fig1e.pdf}\\[-.05cm]
\scriptsize Bulk-boundary
\end{columns}
\end{frame}

\end{spacing}
\end{document}
